\documentclass[11pt]{article}
\usepackage[margin=1in]{geometry}
\usepackage[T1]{fontenc}
\usepackage[utf8]{inputenc}
\usepackage{lmodern,microtype,amsmath,amssymb,amsthm,mathtools}
\usepackage[colorlinks=true,linkcolor=blue,citecolor=blue,urlcolor=blue]{hyperref}
\hypersetup{pdftitle={Superpolynomial sample lower bounds below the mean-field scaling},pdfauthor={Dmitry Vaintrob; Kaarel Hanni; OpenAI Codex}}
\setlength{\emergencystretch}{2em}
\newtheorem{theorem}{Theorem}[section]
\newtheorem{lemma}[theorem]{Lemma}
\newtheorem{proposition}[theorem]{Proposition}
\newtheorem{corollary}[theorem]{Corollary}
\theoremstyle{remark}\newtheorem{remark}[theorem]{Remark}
\newcommand{\E}{\mathbb E}\newcommand{\Pp}{\mathbb P}
\newcommand{\R}{\mathbb R}\newcommand{\X}{\mathcal X}
\DeclareMathOperator{\sign}{sign}\DeclareMathOperator{\Var}{Var}
\title{Superpolynomial sample lower bounds for globally margin-conditioned\ Bayesian neural networks below the mean-field scaling}
\author{Dmitry Vaintrob \and Kaarel H\"anni \and OpenAI Codex}
\date{October 15, 2026}
\begin{document}\maketitle
\begin{abstract}
We study fully connected Bayesian $\tanh$ networks on the Boolean cube,
conditioned before training to have a fixed absolute output margin at
every input. Throughout the positive subcritical scaling regime
$0<\gamma<1$, we prove a superpolynomial sample lower bound for learning
parity. At any fixed depth, after taking the width limit, at least
$\exp(\Omega(\sqrt{k}\log k))$ examples are necessary to reduce the
posterior expected error below $2/5$ on $k$-bit parity. The same sample
lower bound is necessary to give the correct label probability strictly
greater than $2/3$ on more than $9/10$ of the inputs. The bounds hold for
every training set of the indicated size. The proof combines the
actual network's moderate-deviation rate, an explicit spectral estimate
for its saturated kernel, and an elementary polynomial-threshold lower
bound. It does not replace the posterior by a Gaussian posterior and
therefore also covers $\gamma\ge1/2$, where higher-order corrections can
select among tied Gaussian-rate minimizers. An $\exp(\Omega(k))$ bound
at fixed positive test error is not established here.
\end{abstract}

\section{Model and main result}
Fix a dimension $d$, depth $L\ge2$, width $N$, margin $a>0$, and
$0<\gamma<1$. On $\X=\{-1,1\}^d$, let
\[
 h^0(x)=(1,x),\qquad h^s(x)=\tanh(W^s h^{s-1}(x))\ (1\le s<L),
 \qquad z_N(x)=W^L h^{L-1}(x).
\]
All hidden layers have width $N$. The readout is linear, there are no
biases after the augmented input, and the independent centered Gaussian
weights have variances
\begin{equation}\label{eq:weights}
 \Var(W^1_{ij})=N^{1-\gamma}/d,\quad
 \Var(W^s_{ij})=N^{-\gamma}\ (1<s<L),\quad
 \Var(W^L_i)=N^{-\gamma-1}.
\end{equation}
The prior is conditioned on the all-input event
$A_a=\{|z_N(x)|>a\text{ for every }x\in\X\}$.
Training then conditions on exact, noiseless signs
$\sign z_N(x_i)=f(x_i)$. Equivalently, the posterior is the original
weight law conditioned on $A_a$ and these observed signs.

For a sample $S$, write
\[
 p_N(x)=\Pp\{\sign z_N(x)=f(x)\mid A_a,S\},\qquad
 R_N(S,f)=2^{-d}\sum_x[1-p_N(x)],
\]
and use the confidence criterion
\begin{equation}\label{eq:confidence}
 F_N(S,f)=2^{-d}\bigl|\{x:p_N(x)\le2/3\}\bigr|.
\end{equation}
Thus $R_N$ is the expected error of a posterior sample, and $F_N$
counts an input as unsuccessful unless its correct-label probability
is strictly greater than $2/3$. These are different quantities; in
particular this paper does not identify either with the error of the
sign of the posterior latent mean.

Here is an explicit version of the lower bound. The more general
spectral statement is Theorem~\ref{thm:spectral} below.

\begin{theorem}[A sample lower bound throughout $0<\gamma<1$]
\label{thm:main}
Let the target be full parity $f(x)=\prod_{j=1}^k x_j$, with $d=k$.
Set
\[
 c_L=(2/\pi)^{L-1},\qquad
 T_k=\lceil(k+1)\log(4(k+1))\rceil+2,\qquad
 r_k=\lfloor\sqrt{k}/16\rfloor.
\]
Suppose $k\ge4096$ and
\begin{equation}\label{eq:sizeconditions}
 c_L\binom{k}{r_k+1}\ge160(k+1),\qquad
 n\le \frac{c_L}{2560T_k}\binom{k}{r_k+1}.
\end{equation}
For every noiseless training sample with at most $n$ distinct inputs,
and for every fixed $0<\gamma<1$,
\begin{equation}\label{eq:mainbounds}
 \liminf_{N\to\infty}R_N(S,f)\ge\frac25,
 \qquad
 \liminf_{N\to\infty}F_N(S,f)\ge\frac1{10}.
\end{equation}
At each fixed depth, the sample threshold in
\eqref{eq:sizeconditions} is $\exp(\Omega(\sqrt{k}\log k))$.
\end{theorem}

The sample locations can be arbitrary and can include repetitions.
In particular, the conclusion holds for independent uniform training
inputs without an exceptional-sample qualification. Every width limit
is taken at fixed $d,L,\gamma,a,S$; the statement about growth in $k$
concerns these limiting learners. No uniform interchange of the width
and dimension limits is asserted.

Parity is a small-circuit target. Appendix~\ref{sec:paritycircuit}
gives a bounded-weight $\tanh$ realization at every fixed depth. The
lower bound therefore concerns statistical behavior of this specified
prior and conditioning rule, rather than the intrinsic information
needed to identify a parity function.

\section{The kernel rate of the actual network}
Let $M=2^d$, and identify output functions with vectors in $\R^M$.
Define
\begin{equation}\label{eq:kernel}
 H(t)=\frac2\pi\arcsin t,\qquad
 K(x,y)=H^{\circ(L-1)}\!\left(\frac{1+x\cdot y}{d+1}\right).
\end{equation}
The positive definiteness of $K$ is proved in the next section.
For each Boolean function $g$, put
\begin{equation}\label{eq:energy}
 E_K(g)=\min_{g(x)z(x)\ge a\ \forall x}\frac12 z^TK^{-1}z.
\end{equation}
The minimum exists since its feasible set is closed and nonempty and
the objective is coercive. Replacing the closed inequalities by strict
ones gives the same infimum, by scaling a feasible vector by $1+\epsilon$.

\begin{proposition}[Posterior concentration on minimum-energy signs]
\label{prop:rate}
For the actual network~\eqref{eq:weights}, at every fixed
$0<\gamma<1$,
\begin{equation}\label{eq:rate}
 \lim_{N\to\infty}N^{-\gamma}\log
 \Pp\{g(x)z_N(x)>a\ \forall x\}=-E_K(g).
\end{equation}
Consequently, after any fixed noiseless training sample, posterior
mass tends to one on the compatible functions minimizing $E_K$.
No assertion about the relative weights of tied minimizers is needed.
\end{proposition}

\begin{proof}
For hidden layer $s$, let
$Q_N^s(x,y)=N^{-1}\sum_i h_i^s(x)h_i^s(y)$.
We first establish, for each fixed $\epsilon>0$,
\begin{equation}\label{eq:freeze}
 \Pp\{\|Q_N^{L-1}-K\|>\epsilon\}
       \le C_\epsilon e^{-c_\epsilon N}
\end{equation}
for all sufficiently large $N$. Any matrix norm can be used, since
$M$ is fixed. Conditional on the preceding layer, the rows of the
next hidden layer are independent bounded vectors. Hoeffding's
inequality controls all their empirical second moments. Their
preactivations have covariance $N^{1-\gamma}Q_N^{s-1}$; at the first
layer the covariance is
$N^{1-\gamma}(1+x\cdot y)/d$. Because $N^{1-\gamma}\to\infty$,
the conditional second-moment map converges to the Gaussian-sign map
\[
 Q\longmapsto \left[H\!\left(
 \frac{Q(x,y)}{\sqrt{Q(x,x)Q(y,y)}}\right)\right]_{x,y}.
\]
This follows by dominated convergence for a Gaussian vector, since
each scalar variance stays bounded away from zero in a neighborhood
of the limiting covariance. The convergence is uniform in such a
neighborhood: a compactness subsequence and a continuous square root
coupling reduce the assertion to the same dominated-convergence
argument. The limiting map is continuous there. Inducting through
the fixed number of layers, with smaller fixed tolerances at earlier
layers, proves~\eqref{eq:freeze}. The first-layer limit normalizes
the diagonal $(d+1)/d$, giving the kernel in~\eqref{eq:kernel}.

Conditionally on the hidden layers, the readout is exactly Gaussian:
\begin{equation}\label{eq:readout}
 Z_N\mid Q_N^{L-1}\sim\mathcal N(0,N^{-\gamma}Q_N^{L-1}).
\end{equation}
Couple it to $Y_N=N^{-\gamma/2}K^{1/2}\xi$, where
$\xi\sim\mathcal N(0,I_M)$ is independent of the hidden layers.
On a fixed sufficiently small neighborhood of the positive-definite
matrix $K$, the square-root map is Lipschitz. Hence, for $\delta>0$,
\[
 \Pp\{\|Z_N-Y_N\|>\delta\}
 \le C_\epsilon e^{-c_\epsilon N}
 +\Pp\{\|\xi\|>\delta N^{\gamma/2}/(C\epsilon)\}.
\]
First divide logarithms by $N^\gamma$ and let $N\to\infty$;
then let $\epsilon\downarrow0$. The resulting upper limit is
$-\infty$, since $\gamma<1$. The coupled vectors are therefore
exponentially equivalent at speed $N^\gamma$.

The fixed-dimensional Gaussian vectors $Y_N$ have large-deviation
rate $I(z)=z^TK^{-1}z/2$, directly from their densities. Exponential
equivalence transfers the same bounds to $Z_N$. Applying the upper
bound to a closed margin orthant and the lower bound to its interior
gives~\eqref{eq:rate}, since their infima agree. Finally, there are
only finitely many Boolean functions for fixed $d$. Dividing each
compatible orthant probability by their sum proves the posterior
concentration claim.
\end{proof}

This is a rate statement for the network itself. It does not require
relative-error approximation of rare probabilities, nor convergence
in total variation to a Gaussian posterior.

\section{Spectral estimates with explicit dimension dependence}
For $U\subseteq[d]$, let $\chi_U(x)=\prod_{j\in U}x_j$.
The characters diagonalize $K$. Write $\lambda_j$ for the eigenvalue
of the normalized operator $M^{-1}K$ on degree-$j$ characters, and
let $v$ denote the smallest eigenvalue of the matrix $K$ itself.

\begin{lemma}[Spectrum and minimum matrix eigenvalue]
\label{lem:spectrum}
Put $D=d+1$, $c_L=(2/\pi)^{L-1}$, and
$T_d=\lceil D\log(4D)\rceil+2$. Then
\[
 \lambda_0=\lambda_1>\lambda_2=\lambda_3>\cdots>0,
 \qquad v=M\lambda_d,
\]
with a final unpaired value when $d$ is even. Moreover,
\begin{equation}\label{eq:spectralbounds}
 \lambda_0\ge\frac{c_L}{D},\qquad
 v\ge\frac{c_L}{4T_d},\qquad
 \lambda_j\le\binom dj^{-1}\quad(0\le j\le d).
\end{equation}
\end{lemma}

\begin{proof}
Write $F=H^{\circ(L-1)}=\sum_{m\text{ odd}}c_m t^m$.
All these coefficients are positive, sum to one, and satisfy
\begin{equation}\label{eq:coeff}
 c_1=c_L,\qquad c_m\ge c_L/m^2\quad(m\text{ odd}).
\end{equation}
Indeed, the coefficient of $t^{2j+1}$ in $H$ is
$(2/\pi)\binom{2j}{j}/[4^j(2j+1)]$, and
$\binom{2j}{j}\ge4^j/(2j+1)$ because the largest of $2j+1$
binomial coefficients is at least their average. Each additional
composition contributes at least its linear coefficient times the
previous series.

On the full $D$-dimensional cube, the degree-$r$ Fourier eigenvalue
of $(u\cdot w/D)^m$ is the probability $p_m(r)$ that the simple
coordinate-flip walk, started at a fixed vertex, reaches a specified
vertex at distance $r$ after $m$ steps. Expanding the power proves
this: the surviving terms are precisely those whose coordinates
selected an odd number of times are that specified support.
The walk recursion yields
\begin{align*}
 D[p_{m+1}(r)-p_{m+1}(r+2)]
 ={}&r[p_m(r-1)-p_m(r+1)]\\
 &+(D-r-2)[p_m(r+1)-p_m(r+3)]
\end{align*}
for $0\le r\le D-2$, omitting terms with zero boundary coefficients.
Starting at time zero, induction gives monotonicity with distance
within each parity class. Thus the full-cube eigenvalues
$\mu_r=\sum_{m\text{ odd}}c_mp_m(r)$ vanish for even $r$ and decrease
strictly through odd $r$: the positive term $m=r$ distinguishes
$r$ from $r+2$. Restricting the first coordinate to $+1$ gives
$\lambda_r=\mu_r+\mu_{r+1}$. This proves the paired ordering and
positive definiteness. The linear coefficient gives the stated
lower bound on $\lambda_0$. The diagonal $K(x,x)=1$ gives
$\sum_j\binom dj\lambda_j=1$, proving the last inequality.

To bound the smallest eigenvalue from below, let
$m_0=\lceil D\log(4D)\rceil$. The Fourier expansion of the walk
shows that for $m\ge m_0$ and a vertex of matching parity, the
constant and all-coordinate characters contribute $2\cdot2^{-D}$.
The sum of absolute values of the other contributions, divided by
$2^{-D}$, is at most
\[
 \sum_{j=1}^{D-1}\binom Dj|1-2j/D|^m
 \le 2[(1+e^{-2m/D})^D-1]<1.
\]
The last inequality follows from
$D e^{-2m/D}\le1/(16D)$ and $(1+t)^D\le e^{Dt}$.
Consequently $p_m(r)\ge2^{-D}$ whenever $m,r$ are odd and
$m\ge m_0$. Using~\eqref{eq:coeff} and comparing a decreasing
series to its integral gives, for every odd $r$,
\[
 \mu_r\ge 2^{-D}c_L\sum_{\substack{m\ge m_0\\m\text{ odd}}}\frac1{m^2}
 \ge\frac{2^{-D}c_L}{2(m_0+2)}.
\]
Exactly one of $d,d+1$ is odd, so multiplying
$\lambda_d=\mu_d+\mu_{d+1}$ by $M=2^{D-1}$ yields
$v\ge c_L/[4(m_0+2)]$.
\end{proof}

\section{Low-energy outputs cannot approximate parity too well}
We include the elementary polynomial-threshold argument of
Aspnes, Beigel, Furst, and Rudich~\cite{ABFR}; no threshold-complexity
result is assumed beyond the proof below. At zeros a threshold
function may use either fixed sign convention.

\begin{lemma}[Polynomial thresholds and parity]\label{lem:ptf}
If $P$ is a real polynomial of degree at most $r<k$ on
$\{-1,1\}^k$, then
\begin{equation}\label{eq:ptf}
 \Pp_x\{\sign P(x)\ne\chi_{[k]}(x)\}
 \ge b_k(r):=2^{-k}\sum_{j=0}^{\lfloor(k-r-1)/2\rfloor}\binom kj.
\end{equation}
The same bound holds against a parity on $k$ of $d$ coordinates
for a polynomial of total degree at most $r$ on $d$ variables.
Also
\begin{equation}\label{eq:bapprox}
 b_k(r)\ge\frac12-\frac{r+2}{\sqrt{2k}}.
\end{equation}
\end{lemma}

\begin{proof}
A sufficiently small constant perturbation makes $P$ nonzero at all
cube points while preserving the chosen threshold function. Let
$t=\lfloor(k-r-1)/2\rfloor$. If its error set had fewer than
$\sum_{j=0}^t\binom kj$ points, the linear space of multilinear
polynomials of degree at most $t$ would contain a nonzero polynomial
$Q$ vanishing on that set. Then $PQ^2$ is nonzero as a cube function,
has degree at most $r+2t<k$, and has nonnegative product with full
parity everywhere, with a strictly positive product somewhere.
Its inner product with parity is positive, contradicting Fourier
orthogonality. This proves~\eqref{eq:ptf}. For a parity on a subset
of coordinates, condition on all the remaining coordinates and
average the same bound.

For completeness, the largest atom of $\operatorname{Bin}(k,1/2)$
is at most $\sqrt{2/k}$. One can obtain this from
$\binom{2m}{m}/4^m\le1/\sqrt{m+1}$, proved by induction using the
ratio $(2m+1)/(2m+2)$; odd $k$ follows from the adjacent even case.
By symmetry, the binomial tail in~\eqref{eq:ptf} differs from
$1/2$ by at most $(r+2)/2$ times the largest atom. This proves
\eqref{eq:bapprox}.
\end{proof}

\begin{lemma}[A cheap compatible output]\label{lem:competitor}
For any target and any sample with at most $n$ distinct inputs,
the smallest compatible energy satisfies
\begin{equation}\label{eq:competitor}
 E_*\le a^2\left(\lambda_0^{-1}+4n/v\right).
\end{equation}
\end{lemma}
\begin{proof}
Start with the constant vector $a\mathbf1$ and subtract $2a$ at
each observed negative-label input. The result has absolute margin
$a$ everywhere and agrees with every observed label. Write it as
$u+w$, with $u=a\mathbf1$ and $w$ supported on at most $n$ points.
The quadratic norm inequality gives
\[
 \tfrac12(u+w)^TK^{-1}(u+w)
 \le u^TK^{-1}u+w^TK^{-1}w
 \le a^2/\lambda_0+4a^2n/v.
\]
This output is feasible for one compatible closed-margin orthant,
which suffices for its rate energy.
\end{proof}

\begin{theorem}[A spectral sample lower bound]\label{thm:spectral}
Let $f=\chi_U$ with $|U|=k$, and let $S$ have at most $n$ distinct
inputs. For any integer $0\le r<k$, put
\[
 \eta=2\lambda_{r+1}\left(\lambda_0^{-1}+4n/v\right).
\]
Throughout $0<\gamma<1$,
\begin{align}
 \liminf_{N\to\infty}R_N(S,f)&\ge b_k(r)-\eta,\label{eq:riskbound}\\
 \liminf_{N\to\infty}F_N(S,f)&\ge
 \max\{0,[3b_k(r)-3\eta-1]/2\}.\label{eq:confbound}
\end{align}
In particular, with $D=d+1$, $T_d$ as in Lemma~\ref{lem:spectrum},
\begin{equation}\label{eq:etaexplicit}
 \eta\le\frac{2[D+16nT_d]}{c_L\binom d{r+1}}.
\end{equation}
\end{theorem}

\begin{proof}
Take any compatible minimizing sign function $g$, and let $z$
minimize its closed-margin energy. In the normalized Fourier basis,
\[
 2E_* =z^TK^{-1}z
       =\sum_{V\subseteq[d]}\frac{\widehat z(V)^2}{\lambda_{|V|}}.
\]
Let $P(x)=\sum_{|V|\le r}\widehat z(V)\chi_V(x)$. The paired
ordering of the eigenvalues and Parseval's identity give
\[
 \E_x[(z(x)-P(x))^2]\le2E_*\lambda_{r+1}.
\]
Since $|z(x)|\ge a$, a disagreement between $g(x)=\sign z(x)$
and the chosen sign of $P(x)$ implies $|z(x)-P(x)|\ge a$.
Markov's inequality and Lemma~\ref{lem:competitor} bound the
fraction of such disagreements by $\eta$. Lemma~\ref{lem:ptf}
then gives $\operatorname{err}(g,f)\ge b_k(r)-\eta$.

The estimate holds for every compatible minimizing sign, so
Proposition~\ref{prop:rate} proves~\eqref{eq:riskbound} irrespective
of the posterior weights among minimizers. At finite width,
\[
 R_N\le F_N+(1-F_N)/3=(1+2F_N)/3,
\]
which yields~\eqref{eq:confbound}. Substituting the three spectral
bounds from Lemma~\ref{lem:spectrum} proves~\eqref{eq:etaexplicit}.
\end{proof}

\begin{proof}[Proof of Theorem~\ref{thm:main}]
Let $d=k$, $r=r_k$, and $B=\binom{k}{r+1}$. Conditions
\eqref{eq:sizeconditions} imply
\[
 k+1+16nT_k\le c_LB/160+c_LB/160=c_LB/80.
\]
Thus $\eta\le1/40$. As $k\ge4096$,
\[
 b_k(r)\ge\frac12-\frac{1/16+2/64}{\sqrt2}
          >\frac{17}{40}.
\]
Theorem~\ref{thm:spectral} gives $R_N$ lower limit at least
$17/40-1/40=2/5$, and $F_N$ lower limit at least $1/10$.
Finally, $r+1=\Theta(\sqrt{k})$ and
$\binom{k}{r+1}\ge[k/(r+1)]^{r+1}$, so the bound on $n$ is
$\exp(\Omega(\sqrt{k}\log k))$ at fixed $L$.
\end{proof}

\section{What the bound does and does not establish}
The result establishes a fixed-positive-error obstruction for the
actual globally conditioned Bayesian learner. In particular, it is
stronger than a statement that the exact parity function has tiny
prior mass or cannot be the most likely exact function before every
point is observed. The confidence conclusion concerns the marginal
label probabilities in~\eqref{eq:confidence}.

The exponential-in-$k$ sample question remains open. The present
method first upper-bounds the energy of the learner's minimizers and
then bounds what any function of that energy can accomplish. Its
limitation is concrete: a polynomial with roots between consecutive
Hamming levels in a central band of width $O(\sqrt{k})$ has degree
$O(\sqrt{k})$, and its sign agrees with parity on a fixed large
fraction of the cube. Scaling that polynomial to a global margin
gives kernel energy $\exp(O(\sqrt{k}\log k))$. Thus energy alone
cannot give an $\exp(\Omega(k))$ necessary scale at fixed error.
A stronger sample lower bound would have to use more about the
compatible minimum-energy interpolants and the training constraints.
Conversely, that approximation does not prove a learning upper
bound: its wrong-sign outputs outside the central band can be very
large, so correcting observations there need not be cheap.

The regime $\gamma\le0$ has a different limiting posterior rather
than the energy-minimization limit used here. The theorem makes no
sample-complexity assertion for that regime.

\appendix
\section{A small bounded-weight parity network}\label{sec:paritycircuit}
For completeness, full parity on $k$ bits has a one-hidden-layer
$\tanh$ realization of width $k+1$, with weights $O(k\log(k+1))$
and output margin greater than $1/2$. Let
\[
 t(x)=\frac{k+\sum_i x_i}{2}\in\{0,\ldots,k\},\quad
 y_j=(-1)^{k-j},\quad A=\log(16(k+1)).
\]
The step functions $s_j(t)=\sign(t-j+1/2)$ satisfy
\[
 y_t=c+\sum_{j=1}^k y_js_j(t),\qquad c=(y_0+y_k)/2\in\{0,1\}.
\]
Replace each step by $h_j=\tanh(A(t-j+1/2))$. Since its argument
has absolute value at least $A/2$, its error is at most $2e^{-A}$.
A constant neuron $h_0=\tanh1$ supplies the term $c$ using output
weight $c/\tanh1$. The resulting output differs from parity by
less than $1/8$. First-layer input weights are $A/2$ and bias
weights have magnitude at most $Ak/2$; readout weights have
magnitude at most $2$.

At depth $L\ge3$, the next hidden layer sums these features and
applies $\tanh$, producing the correct sign with magnitude at least
$\tanh(7/8)>1/2$. Each additional hidden layer can preserve this
margin by a single weight $2$, since $\tanh(2u)>1/2$ whenever
$u>1/2$. A final readout weight $1$ completes the construction.
Unused neurons are set to zero. Thus the same polynomial bounds
hold at every fixed depth.

\begin{thebibliography}{9}
\bibitem{ABFR} J. Aspnes, R. Beigel, M. Furst, and S. Rudich.
\emph{The expressive power of voting polynomials}.
Combinatorica 14(2):135--148, 1994.
\href{https://www.cs.yale.edu/homes/aspnes/papers/stoc91voting.pdf}{Author-hosted manuscript}.
\end{thebibliography}
\end{document}
